% Generated draft: requires mathematical and source review.
\documentclass[11pt]{article}
\usepackage[margin=25mm]{geometry}
\usepackage{fontspec}
\setmainfont{Latin Modern Roman}
\usepackage{amsmath,amssymb,mathtools,unicode-math}
\setmathfont{Latin Modern Math}
\usepackage{enumitem,needspace,xurl,hyperref}
\setlength{\parindent}{0pt}
\setlength{\parskip}{5pt}
\setlength{\emergencystretch}{3em}
\newcommand{\N}{\mathbb{N}}
\newcommand{\Z}{\mathbb{Z}}
\newcommand{\Q}{\mathbb{Q}}
\newcommand{\R}{\mathbb{R}}
\newcommand{\C}{\mathbb{C}}
\newcommand{\F}{\mathbb{F}}
\newcommand{\A}{\mathbb{A}}
\newcommand{\PP}{\mathbb{P}}
\newcommand{\E}{\mathbb{E}}
\newcommand{\eps}{\varepsilon}
\newcommand{\dd}{\,\mathrm{d}}
\begin{document}
\section*{UnsolvedMath ID: 30001894}
\textbf{Source problem code:} catalogue entry
\subsection{Definitions and mathematical statement}
Let \(K\subseteq\R^2\) be a compact convex set, and let \(\mathcal T(K)=\{K+t:t\in\R^2\}\), where \(K+t=\{x+t:x\in K\}\). Let \(F_1,F_2,F_3\subseteq\mathcal T(K)\). Under the point-transversal interpretation, a \(3\)-transversal of a family \(F\) is a set \(P\subseteq\R^2\) with \(|P|\leq 3\) such that \(P\cap A\neq\varnothing\) for every \(A\in F\). Write this property as
\[
T_3(F)\iff\exists P\subseteq\R^2\;\bigl(|P|\leq3\ \land\ \forall A\in F,\ P\cap A\neq\varnothing\bigr).
\]
The supplied question asks whether the stated pairwise-intersection hypothesis forces
\[
\exists i\in\{1,2,3\}\;T_3(F_i).
\]
However, the supplied record does not define either ``pairwise intersecting families'' or ``\(3\)-transversal.'' Two distinct possible intersection hypotheses are
\[
\begin{aligned}
H_{\mathrm{cross}}&:\quad
\forall i\neq j\;\forall A\in F_i\;\forall B\in F_j,
\quad A\cap B\neq\varnothing,\\
H_{\mathrm{within}}&:\quad
\forall i\in\{1,2,3\}\;\forall A,B\in F_i,
\quad A\neq B\Longrightarrow A\cap B\neq\varnothing.
\end{aligned}
\]
For each choice \(H\) among these hypotheses, a precise candidate question is whether, for every compact convex \(K\subseteq\R^2\) and every three families \(F_1,F_2,F_3\subseteq\mathcal T(K)\),
\[
H\Longrightarrow\exists i\in\{1,2,3\}\;T_3(F_i).
\]
These are alternative formulations requiring source clarification, not an assertion that either is the intended formulation. The point-transversal interpretation also requires confirmation. The record does not specify whether the families must be finite, whether \(K\) must be nonempty or have nonempty interior, or whether empty families are permitted; the intended quantified domain therefore remains unresolved.
\subsection{Short English statement}
The supplied question asks whether three suitably intersecting families of translates of one compact convex planar set must include a family that can be met by three points, under the point-transversal interpretation. The meanings of the intersection condition and transversal, and the permitted sets and families, require clarification before a unique precise statement can be selected.
\subsection{Sources}
\begin{itemize}
\item Local UnsolvedMath registry, supplied complete source record, permanent ID \(30001894\), title ``Colored Transversals for Translates of Convex Sets,'' source problem code \texttt{catalogue entry}. The project registry locations are \texttt{data/unsolvedmath/problems.json} and \texttt{catalog/problems.json}; no independent file inspection was performed.
\item Supplied report citation: \emph{Discrete Geometry (2012)}, Oberwolfach Reports, DOI 10.4171/owr/2011/44, pp.\ 2459--2548. \url{https://doi.org/10.4171/owr/2011/44}. The literature-source entry labels this same reference ``OWR 44 (2012).'' The supplied proposed year is 2011; no individual proposer is identified.
\item Literature assessment dated exactly \texttt{2026-08-21}: the supplied curation labels the problem open and reports that no general resolution of the colored Grunbaum transversal question was found, and that the report presents it as open. This is supplied dated curation, not an independently verified literature assessment or verification of the cited report. Its listed checked source is \url{https://doi.org/10.4171/owr/2011/44}.
\item The supplied statement status is \texttt{corrected\_verified}; its accompanying explanation says that the complete statement was preserved while mathematical notation was normalized. This supplied label does not establish independent mathematical verification of the formulation above.
\end{itemize}
\end{document}
